\section{Controlled Tests on the Sample}
\label{sec:sample}

This section tests the theory where it can be checked most precisely: on the
5,000-node sample and its 200-query benchmark (Section~\ref{sec:data-sample}),
with thousands of independent damage orders (Figure~\ref{fig:sample}).

\begin{figure}[!htbp]
  \centering
  \includegraphics[width=0.72\linewidth]{act0/fig_sample_curves.pdf}
  \caption{The controlled benchmark on the sample under random edge removal:
  giant component $S$ (mean of ten damage orders) and retrieval success $R_L$
  (mean of 600 independent damage orders). Shading shows the spread of a single
  damage state (one standard deviation). Dashed lines are the formulas of
  Section~\ref{sec:theory}, which contain no fitted parameters.}
  \label{fig:sample}
\end{figure}

\subsection{Connectivity and two-hop retrieval}
\label{sec:blind-connectivity}

Under random edge removal the giant component barely shrinks until most edges
are gone: the ten-order means of $S$ are 0.982 at $f=0.5$, 0.815 at $f=0.9$ and
0.112 at $f=0.99$. The generating-function predictions,
Eqs.~\eqref{eq:gf-bond}--\eqref{eq:gf-site}, stay within 0.0011 (edges) and
0.0033 (nodes) of these means for $f\le0.9$, which mainly confirms the local
limit, Eq.~\eqref{eq:s-local}: a node drops out when it loses its last edge.
The full graph, with lower mean degree, holds 0.858 and 0.554 of its nodes at
$f=0.5$ and $0.9$, for the same reason.

Two-hop retrieval follows Eq.~\eqref{eq:r2} to Monte Carlo precision: over
5,000 independent orders, the largest difference at any damage level is 0.0008
for edge removal and 0.0010 for node removal, at most 2.4 standard errors. This
confirms the code and the independence of routes. The shaded bands are much
wider than this agreement because they show a single damage state, which
scatters by 0.047 at $f=0.5$: 2.2 times the variance expected for independent
queries, because the 200 queries share 39 drugs and behave like about 90
independent ones. Fixed damage orders are nested prefixes of one permutation,
so their deviations at different $f$ are correlated (0.58) and are not
independent evidence. Redundancy matters a
great deal at moderate damage and little at heavy damage: at $f=0.3$ a second
shared protein raises survival from 0.49 to 0.74, but at $f=0.7$ only from 0.09
to 0.17. Because Eq.~\eqref{eq:r2} needs only the redundancy distribution, it
also extends beyond the benchmark: among the 406,600 drug--disease pairs of the
full graph that share a protein, 81\% have $m=1$, and only 30\% keep a two-hop
link at $f=0.5$.

\subsection{The horizon}
\label{sec:blind-horizon}

At $f=0.5$ the long-run means are $S=0.982$, $R_\infty=0.862$, $R_4=0.861$,
$R_3=0.792$ and $R_2=0.414$. The 57-point gap between connectivity and two-hop
retrieval therefore has two parts. About 12 points come from the step between
counting nodes and counting queries: a query needs both endpoints attached, and
its drug often has few links. The remaining 45 points come from the hop limit.
On the full graph the first part vanishes (91\% of two-hop indications remain
reachable at $f=0.5$, more than the 86\% of nodes in the giant component), so the
gap there is entirely due to the hop limit.

For unlimited hops, the attachment approximation of Section~\ref{sec:theory-rinf} agrees with the long-run means to within
0.009 at every damage level from $f=0.05$ to $0.99$ (Figure~\ref{fig:sample}).
The simpler endpoint-degree form, Eq.~\eqref{eq:mincut}, agrees to within 0.005 for
$f\le0.8$. Beyond that it overestimates success, by 0.021 at $f=0.9$, because
the PPI giant component begins to shrink. Even so, the PPI layer outlasts the
query endpoints. Its giant component still holds 75\% of proteins at $f=0.9$
(generating-function value 0.751), and its percolation threshold, Eq.~\eqref{eq:pc}, lies beyond
99\% edge loss. Up to $f=0.9$, more than 90\% of the queries that fail at
$L=\infty$ have lost every link of their drug while the disease side is still
attached. Longer horizons therefore move the bottleneck from shared proteins to
drug target links.

This bottleneck depends on the drug degree distribution. In the sample, 69\% of
drugs have one target. In the full filtered graph the figure is 53\%, so drug
links are a weak point of PrimeKG as a whole, but the size of the effect here
is specific to the sample.


Limited-path percolation, Eq.~\eqref{eq:lpp}, is a negative result here
(Table~\ref{tab:lpp}). Its threshold assumes that the number of paths grows by
$\kappa_0-1\approx188$ per extra step. On the sample it matches the damage at
which $R_L/R_\infty=1/2$ to within 0.016, but the sample's queries have a median
ratio $n_3/n_2$ of 25, not 188, as the chance-path calculation of
Section~\ref{sec:theory-chance} also suggests for the first step (24 on the
full graph), so this agreement is most likely a coincidence. On the full graph the threshold fails: the
crossings for two-hop indications come at $f=0.63$ and $0.85$ against 0.83 and
0.93, and two-hop decoys behave the same way (0.61 and 0.84). The local estimate,
Eq.~\eqref{eq:paths}, which uses each pair's actual path counts, follows the
full graph's $R_3$ to within 0.06 at the four damage levels we tested, and
overestimates the sample's by up to 0.13, because there many three-hop paths
share the drug's only edge. Pairs in the full graph have few longer paths (a
median of 5 three-hop paths per two-hop indication, against 37.5 per benchmark
query), so longer horizons protect them less than a mean-field threshold
suggests.

\begin{table}[!htbp]
\centering
\small
\caption{Limited-path percolation. Predicted and measured damage $f$ at which
$R_L/R_\infty=1/2$, and the median numbers of two- and three-hop paths per query
in the intact graph. Sample: 600 orders. Full graph: five orders.}
\label{tab:lpp}
\begin{tabular}{@{}lccccc@{}}
\toprule
Query set & $L$ & Predicted, Eq.~\eqref{eq:lpp} & Measured & Median $n_2$ & Median $n_3$ \\
\midrule
Sample benchmark          & 3 & 0.825 & 0.810 & 1.5 & 37.5 \\
                          & 4 & 0.927 & 0.921 &     &      \\
Full graph, two-hop indications & 3 & 0.826 & 0.629 & 1   & 5    \\
                          & 4 & 0.927 & 0.846 &     &      \\
\bottomrule
\end{tabular}
\end{table}

\subsection{Null models for the excess fragility}
\label{sec:targeted-null}

The degree attack fragments the sample faster than Callaway's theory predicts:
at $f=0.3$ the theory gives $S=0.648$ where the graph has 0.546, while the
theory is accurate to 0.0012 for stub-matched multigraphs with the same degrees.
The excess must come from the requirement that the real graph be simple, or
from structure that the degree sequence does not record. To attribute it, we
rewired the sample in three ways, each keeping every node's degree so that the
same ranking applies: full rewiring \cite{maslov2002specificity} keeps only the
degree sequence; typed rewiring swaps edges within one relation, so every node
keeps its degree in each relation; and 2K rewiring swaps only edge ends of equal
degree, keeping the joint degree distribution and hence all degree--degree
correlations. Each model was realized five times with ten attempted swaps per
edge; afterwards only 9\%, 13\% and 14\% of the original edges remain, and the
standard deviation of $S$ across realizations is at most 0.015 except near the
collapse.

\begin{figure}[!htbp]
  \centering
  \includegraphics[width=0.9\linewidth]{act1/primekg_5000_bfs_v1/fig_attack_null_model.pdf}
  \caption{Degree attack on the sample and on three degree-preserving
  null models. (a) Giant component. (b) Difference from the Callaway prediction.
  Null models are means over five realizations. Shading shows one standard
  deviation.}
  \label{fig:null}
\end{figure}

For $f\le0.5$, Figure~\ref{fig:null} orders the models as expected. Full
rewiring is closest to the theory, typed rewiring is lower, 2K rewiring is lower
still, and the real graph is lowest. We split the gap between theory and the real graph into three
steps. The step from theory to full rewiring is the cost of requiring a simple
graph with a heavy-tailed degree sequence. The step from full to 2K rewiring is
the effect of degree correlations beyond those that the simple-graph constraint
already forces (fully rewired graphs have $r=-0.108$). The step from 2K rewiring to the real graph
is everything else. For $0.1\le f\le0.4$, these steps account for 6--18\%,
32--50\% and 44--53\% of the gap. As $f$ grows, the share of degree correlations
falls and the share of the simple-graph constraint rises.

The mechanism behind the degree-correlation step is visible directly. In the
sample, low-degree nodes hang from hubs (Table~\ref{tab:graphs}), so removing
the hubs strands them. A 5\% attack isolates 6.7\% of all nodes in the real
graph, 4.3\% in the 2K graphs and 1.1\% after full rewiring. Drugs are the
extreme case. After the attack, 77\% of the surviving drugs have no neighbor
left in the real graph and 78\% in the 2K graphs, but only 26\% after full
rewiring. The real graph is more clustered than
the 2K graphs (transitivity 0.069 against 0.051). Clustering is therefore one
candidate for the residual half, together with mesoscale structure such as
disease modules and the BFS sampling itself. We did not separate these.

To separate PrimeKG from its sample, we repeated the degree attack on the full
filtered graph. It also falls below its Callaway prediction, by up to 0.06 near
$f=0.25$ (0.352 against 0.415), before both collapse near $f\approx0.35$. At equal
$f$, the sample's absolute gap is only 1.3--1.4 times larger (0.081 against
0.062 at $f=0.2$). Relative to the predicted value, the full graph is as fragile
or more so (a 12\% shortfall against 10.5\% at $f=0.2$). The excess fragility
under attack is therefore a property of PrimeKG, not an artefact of the sample,
even though the two graphs differ in their degree correlations.

\subsection{The benchmark under attack}
\label{sec:hubs-sample}

Removing 5\% of nodes leaves 87.9\% of the benchmark queries according to
Eq.~\eqref{eq:r2-node} and 87.0\% after random removal, but 30.5\% after the
degree attack and 24.0\% after the betweenness attack. As on the full graph,
the 69.5 points lost to the degree attack split almost evenly between lost
endpoints, mostly hub diseases (34.0), and lost routes (35.5). The giant
component again understates the damage, by a factor of about 8.5, and it even
ranks the two attacks the wrong way round: by $S$ the degree attack is
marginally worse (0.883 against 0.886), by $R_2$ the betweenness attack is worse
by 6.5 points.
